\section{Source audit, reproducibility, and further research}
\label{xid:sec:audit}

\subsection{Corrections and exact claim status}

Two narrow issues remain in the pinned canonical
\path{chapters/03-algebraic.tex}. Both had been flagged in an earlier
continuation and are confirmed here; they are not presented as newly
discovered mistakes. The proposed patch preserves the existing labels
and the proved identities.

\paragraph{The real logarithm in \texttt{cleo:eq:lintanh}.}
The formula concerns real $|\lambda|<1$, with
$\theta=\arcsin\lambda$. Its term
$\theta\log\tan(\theta/2)$ must be
\[
 \theta\log\left|\tan\frac\theta2\right|,
\]
with continuous value zero at $\theta=0$, unless the parameter range is
restricted to positive $\lambda$. For negative $\lambda$ the displayed
logarithm otherwise introduces an imaginary term into a real integral.
For example, at $\lambda=-1/2$ the spurious principal-logarithm term is
$-i\pi^2/6$. The surrounding Clausen terms are unchanged. The real
identity is checked by differentiation in $\theta$, giving
$\theta/\sin\theta$ on both sides, and the zero value at the origin.

\paragraph{The description of the unweighted diagonal.}
The phrase ``the simplest non-elementary diagonal'' is replaced by
``the $(1,1)$ diagonal has the Legendre--chi representation.'' The exact
formula
\[
 \int_0^{\pi/2}\arctan^2(\sin x)\,dx
       =\pi\chi_2(3-2\sqrt2)
\]
is retained. A representation alone does not prove minimality or
nonreduction in a specified algebra of constants. This correction does
not concern the different, weighted integral evaluated in
\eqref{xid:wat:eq:K11}.

\medskip
\noindent The package contains a unified patch and its before/after
hashes. The patch passed a dry-run application and was applied to a
temporary copy, whose bytes matched the intended correction. The
baseline itself was unchanged.

\begin{longtable}{@{}>{\raggedright\arraybackslash}p{0.29\textwidth}>{\raggedright\arraybackslash}p{0.64\textwidth}@{}}
\toprule
Source item & Status after this report\\
\midrule
\endfirsthead
\toprule Source item & Status after this report\\\midrule
\endhead
\texttt{gauss:eq:S4-closed} & Already proved in the canonical
\path{04-S4-proof.tex}; this report does not claim a new proof of its
rational octahedral certificate.\\[3pt]
\texttt{cycloquot:conj:S6} & The short $S_6$ identity remains
conjectural. The existing incomplete-beta and reflection generators do
not supply the missing odd-weight relation.\\[3pt]
\texttt{s8new:conj:S8} & The current $S_8$ candidate remains
conjectural. It is distinct from the rejected earlier vector recorded
under \texttt{research:prop:S8-rejected}.\\[3pt]
Higher golden vectors & The numerical $\Li_6$--$\Li_9$ candidates in
\path{03-algebraic.tex} have not acquired proofs here.\\[3pt]
Colliding unequal grids & Resolved for every pair of Stieltjes indices
and argument orders by Section~\ref{xid:cld:section}, including the full
subtraction polynomial and convergent collision matching.\\[3pt]
Nonlinear local coordinates & Resolved for orientation-preserving
analytic local diffeomorphisms by Section~\ref{xid:sec:nonlinear}, including
the composition law and finite-jet dependence.\\[3pt]
Three unequal frequencies & Resolved under pairwise grid separation
by Section~\ref{xid:sec:triple-dilation}, including all argument and order
coefficients. Colliding triples remain a separate problem.\\[3pt]
Fully shifted height-one sums & An all-depth extension, with the
depth-one literature credited. The inner shift must not be confused
with the older outer-only shift or with simultaneous changes of all
exponents.\\[3pt]
Weighted arctangent family & Evaluated for all real parameters, with
ordinary and Euler primitives and a direct inverse-hyperbolic companion.
No global priority or minimality claim is made.\\
\bottomrule
\end{longtable}

The separated-grid theorem in the incoming report is valid under its
stated separation hypothesis. Extending its distributional product
directly through collision would be an invalid use of that theorem,
rather than an error in the separated statement. Similarly, omitting
$\Delta_{m,r}(a)$ from the fully shifted harmonic formula would give an
incorrect extension; the present report does not attribute that omitted
correction to an existing source. These distinctions matter when the new
sections are integrated.

\subsection{What the computational artifacts establish}

The analytic proofs establish the parameter-uniform identities. The
computer checks target signs, scales, coordinate constants, finite
coefficient extraction, and independent numerical evaluations. They do
not infer identities from small residuals. All numerical results are
ordinary floating-point audits; no result file is described as an
interval certificate.

\begin{center}
\small
\begin{tabular}{@{}>{\raggedright\arraybackslash}p{0.23\textwidth}>{\raggedright\arraybackslash}p{0.42\textwidth}>{\raggedright\arraybackslash}p{0.25\textwidth}@{}}
\toprule
Family & Independent audit & Observed residual\\
\midrule
Nonlinear coordinates & 213 exact rational-series checks, including
cutoff constants and composition; 4 subtracted integrals at 70 digits
& Absolute $<4.0\cdot10^{-36}$\\[4pt]
Colliding products & 33 exact coefficient checks; 9 polygamma integrals
at 44 digits and 2 Stieltjes derivative integrals at 36 digits
& Scaled $<6.7\cdot10^{-45}$ and $<2.8\cdot10^{-31}$\\[4pt]
Fully shifted sums & 72 checks at 75 digits, using Gamma-ratio
Dirichlet tails and direct local differentiation
& Absolute $<8.5\cdot10^{-52}$\\[4pt]
Three frequencies & 2 exact rational Bernoulli integrals compared with
12 weighted-kernel Mellin integrals at 60 digits
& Absolute $<7.2\cdot10^{-66}$\\[4pt]
Weighted primitives & 78 comparisons at 80 digits; symbolic checks of
the mixed derivatives, primitive, and coefficient obstruction
& Absolute $<1.3\cdot10^{-79}$\\
\bottomrule
\end{tabular}
\end{center}

There are 167 numerical comparisons in total. The difference between
working precision and observed residual is expected: some checks use
finite Taylor or asymptotic tails, while others compare small exact
values with direct quadrature. The results record the actual precision
and method rather than a uniform advertised number of correct digits.
For the harmonic continuation, two larger-tail repeats give residuals
below $1.7\cdot10^{-64}$, which tests the identified truncation risk.

The seven mathematical scripts are independent of the TeX build. They
use Python, SymPy, and mpmath. A driver reruns them, writes their logs,
and validates the result files. The exact collision generator also
accepts $m,n,r,k$ as command-line arguments and prints the corresponding
finite symbolic identity, with formal symbols for ordinary zeta
derivatives. The package includes the source manifest, a claim ledger,
integration instructions, the correction patch, and both modular and
standalone TeX. No numerical relation search is needed to reproduce any
proof in this report.

\subsection{Proposed exact research problems}
\label{xid:sec:further-research}

The following questions concern identities and structural reductions.
They are questions, not additional theorems. Each starts from a specific
formula proved above.

\begin{enumerate}[label=\textbf{Q\arabic*.},leftmargin=*]
\item \textbf{Mixed collisions of three unequal grids.}
Remove pairwise separation in Theorem~\ref{xid:thm:triple-FP}. Classify
pairwise collision points and triple collision points, write their
complete local subtraction polynomials, and determine whether the
regular constant in every simultaneous collision limit agrees with the
merged ambient finite part. The two-grid result suggests a positive
matching statement for suitably specified collision coordinates, but
multiple rates can introduce distinct logarithmic terms and must be
tracked explicitly.

\item \textbf{An all-factor finite congruence formula.}
For $\ell\ge4$ pairwise separated integer grids, generalize the six
sign-cone proof to a finite collection of colored multiple Tornheim or
Mordell--Tornheim kernels on the rank-$(\ell-1)$ lattice
$\sum p_jn_j=0$. Determine a useful minimal finite character formula,
then give every ambient-coordinate Stieltjes coefficient with complete
local numerators. The challenge is an exact finite organization of the
lattice, rather than existence of meromorphic continuation alone.

\item \textbf{Which odd-order cancellations survive at more than two
factors?}
Corollary~\ref{xid:cld:diagonal-collapse} cancels an entire antisymmetric
two-factor spectral kernel. Identify a permutation or character
projection for three or more factors that removes the genuinely
multiple spectral terms and leaves only local Stieltjes data. A useful
first case has equal Stieltjes indices and one total argument derivative.

\item \textbf{Nonlinear transport of merged products.}
Theorem~\ref{xid:thm:coordinate-monomial} applies to the full singular
expansion of a colliding product. Compose it with
\eqref{xid:cld:common-subtraction} to obtain a single closed multivariate
generator for every nonlinear pullback of every two-grid product.
Determine sharp vanishing thresholds in the two Stieltjes indices.
For a nonunit tangent, factor out the affine contribution before
classifying the curvature terms.

\item \textbf{A finite algebra for coordinate composition.}
Express the coefficients of \eqref{xid:eq:C-generator} and
\eqref{xid:eq:delta-pullback} in a consistently normalized Bell-polynomial
basis and identify the resulting composition matrices. Which
logarithmic singular germs have zero defect for every coordinate tangent
to the identity? The theorem supplies a direct coefficient test, but a
classification by singular order and logarithmic degree could expose
additional exact cancellations.

\item \textbf{Higher regular Laurent coefficients of the harmonic
entire difference.}
Evaluate $\partial_s^kD(-m,u;a)$ for $k\ge1$. Its value at $k=0$
is a finite Stirling sum, whereas differentiation introduces logarithmic
Mellin moments of the analytic Gauss block. Seek closed expressions in
Gamma derivatives, generalized Stieltjes constants, and a precisely
specified class of multiple polylogarithms. Do not assume the rational
polynomial closure at $k=0$ persists for $k\ge1$.

\item \textbf{Identify the Stirling polynomial system exactly.}
Compare
$\sum_j\left\{\begin{smallmatrix}m+1\\j\end{smallmatrix}\right\}
(a-1)^{\underline j}/j^{r+1}$
with poly-Bernoulli and related polynomial systems, including the
falling-factorial convention and the shift by one. An exact generating
function identification could yield addition, reflection, and
differentiation identities for the complete harmonic correction.

\item \textbf{Independent inner exponents.}
Replace the inner exponent string $(1,\ldots,1)$ in $M_r$ by independent
parameters, keeping the common shift. Determine the smallest generating
kernel whose singular part still separates by a Gamma factor, and
identify the extra entire corrections. This would connect the present
height-one calculus to the incoming all-exponent nested jets without
conflating their spectral directions.

\item \textbf{Rational shifts and character projections.}
Apply finite Dirichlet-character sums to the $a$-dependence of the
Laurent formulas. Determine exact reductions of their Gamma-polygamma
coefficients to Dirichlet $L$-values and their derivatives, and identify
which character parities eliminate the polynomial Stirling correction.
The finite correction should be summed before imposing parity.

\item \textbf{A real addition law for higher Euler primitives.}
The first weighted arctangent integral has the two-term formula
\eqref{xid:wat:eq:main}; the higher Euler primitives currently have the
finite depth-two expression \eqref{xid:wat:eq:Imclosed}. Determine which
higher integration orders admit a comparable addition formula in
$\alpha\pm\beta$, and which require irreducible two-variable functions.
The known exact coefficient obstruction concerns one product variable
only; it does not rule out these other organizations.

\item \textbf{Relations among the algebraic endpoint values.}
The weighted Cleo identity combines logarithms of the silver and golden
units with $\chi_2$ and $\chi_3$ at explicit algebraic arguments. Seek
functional-equation certificates for further reductions at these
arguments. A target basis should be specified in advance, and failure
of a finite relation search should not be interpreted as independence.

\item \textbf{The short $S_6$ and current $S_8$ candidates.}
Extend the convergent octahedral or lifted functional-equation mechanism
that proved $S_4$ to weights seven and nine. The existing reflection
projection cancels at the relevant odd total weights, so a successful
proof needs an additional functional relation. A useful deliverable
would be an exact rational row certificate with an independently
checked analytic origin for every row. The frozen candidate vectors,
not newly fitted alternatives, should remain the targets.

\item \textbf{Formalization of the local finite-part lemmas.}
Formalize the elementary Taylor-subtraction coefficient lemma first,
including its exact test-function and coordinate conventions. It can
then serve as a common foundation for nonlinear contacts and colliding
products. The finite rational coefficient computations are suitable
for checked algebraic certificates, while the analytic continuation
and dominated-differentiation steps remain separate proof obligations.
\end{enumerate}

The most immediate continuations are Q1, Q4, and Q6: their inputs are
explicitly available in this report, and each asks for an exact identity
whose first unresolved term can be computed without a numerical basis
guess. The $S_6$ and $S_8$ problems remain more structurally demanding;
the present endpoint calculus supplies normalization discipline, not
the missing weight-seven or weight-nine functional equation.
